% Work-MAT, hodina 58. Zaverecne opakovani celeho bloku algebraickych vyrazu 41--57.
% Podklady: ucebnice Algebraické výrazy, metodicka inspirace realisticky.cz a
% puvodni priprava 58 v repozitari. Nezavadi se nove ucivo.
\documentclass[11pt,a4paper]{article}
% Spolecny styl pro ucitelske pripravy MAT (Work/MAT)
% Verze 1.1, 2026-10-23
\NeedsTeXFormat{LaTeX2e}
\ProvidesPackage{priprava}[2026/10/23 v1.1 Sablona ucitelskych priprav MAT]

% Zakladni typografie a cestina
\RequirePackage[T1]{fontenc}
\RequirePackage[utf8]{inputenc}
\RequirePackage[czech]{babel}
\RequirePackage{lmodern}
\RequirePackage{microtype}
\RequirePackage{amsmath,amssymb}
\RequirePackage{xcolor}
\RequirePackage{enumitem}
\RequirePackage[most]{tcolorbox}
\RequirePackage{tabularx,booktabs,array}
\RequirePackage{fancyhdr}
\RequirePackage{needspace}
\RequirePackage{geometry}
\RequirePackage[hidelinks]{hyperref}

\geometry{a4paper,top=17mm,bottom=18mm,left=23mm,right=23mm}
\definecolor{notegray}{gray}{0.39}
\setlength{\parindent}{0pt}
\setlength{\parskip}{2.5pt plus 1pt minus 0.5pt}
\setlist[itemize]{topsep=2pt,itemsep=1pt,parsep=0pt,leftmargin=1.45em}
\setlist[enumerate]{topsep=2pt,itemsep=2pt,parsep=0pt,leftmargin=1.65em}
\setcounter{secnumdepth}{0}
\renewcommand{\arraystretch}{1.12}

% Zahlavi: vlevo rocnik | globalni tema, vpravo cislo: nazev hodiny
% Zapati zustava zcela prazdne.
\pagestyle{fancy}
\fancyhf{}
\renewcommand{\headrulewidth}{0.35pt}
\renewcommand{\footrulewidth}{0pt}
\setlength{\headheight}{14pt}
\newcommand{\nastavPripravu}[4]{%
  \fancyhead[L]{\small #1\enspace\textbar\enspace #2}%
  \fancyhead[R]{\small #3: #4}%
}

% Minimalisticke cernobile boxy. Kazda definice zvlast.
\tcbset{commonlesson/.style={enhanced,breakable=false,colback=white,colframe=black,
  colbacktitle=white,coltitle=black,boxrule=0.45pt,
  arc=0.8mm,outer arc=0.8mm,left=3mm,right=3mm,top=1.6mm,bottom=1.6mm,
  toptitle=1mm,bottomtitle=0.7mm,fonttitle=\bfseries,
  attach boxed title to top left={xshift=2.5mm,yshift=-1.15mm},
  boxed title style={colback=white,colframe=white,boxrule=0pt,left=1mm,right=1mm}}}
\newtcolorbox{zadaniBox}[1][]{commonlesson,title={Zadání},#1}
\newtcolorbox{prikladBox}[1][]{commonlesson,title={Příklad},#1}
\newtcolorbox{definiceBox}[1][]{commonlesson,title={Definice},#1}

% Spolecny box pro zadani a jeho reseni, bez mezery mezi castmi.
% Pouziti:
%   \begin{prikladReseniBox}
%     Zadani ulohy.
%     \tcblower
%     \begin{align*} ... \end{align*}
%   \end{prikladReseniBox}
% Spodni cast ma jemny sedy podklad a delici linku; lze v ni sazet
% vice radku matematiky. Obe casti maji jen jeden spolecny ramecek.
\tcbset{solutionlesson/.style={
  commonlesson,
  enhanced jigsaw,
  breakable,
  colbacklower=black!3,
  segmentation style={solid,draw=black!25,line width=0.35pt},
  fontlower=\small\color{notegray},
  before lower={\noindent\textbf{Řešení:}\enspace},
  leftlower=3mm,rightlower=3mm
}}
\newtcolorbox{prikladReseniBox}[1][]{solutionlesson,title={Příklad},#1}
\newtcolorbox{zadaniReseniBox}[1][]{solutionlesson,title={Zadání},#1}

% Pomucky pro autora pripravy: sedive poznamky; \reseni pro starsi pripravy
\newcommand{\poznamka}[1]{\par\begingroup\small\color{notegray}\itshape
  \textnormal{\textbf{Poznámka pro učitele.}}\enspace #1\par\endgroup}
\newcommand{\reseni}[1]{\par\begingroup\small\color{notegray}
  \textbf{Očekávané řešení:}\enspace #1\par\endgroup}
\newcommand{\blok}[1]{\needspace{5\baselineskip}\subsection*{#1}}
\newcommand{\vyrok}[1]{\uv{#1}}

\nastavPripravu{1. ročník}{Algebraické výrazy}{58}{Opakování a ověřování správnosti úprav}
\begin{document}
\poznamka{45 minut: 5 min strategie kontroly, 10 min společná práce s identitou, 15 min samostatná práce a kontrola, 10 min práce s chybným výpočtem, 5 min shrnutí. Rozšíření jen pro rychlé, domácí úlohy 15--20 minut. Navazuje na diagnostiku hodiny 57.}
\blok{Kdy máme jistotu, že jsme výraz upravili správně?}
Dosazení několika čísel je \emph{kontrolou}; neshoda v jediném přípustném dosazení rovnost vyvrací. K obecné platnosti potřebujeme zdůvodněné algebraické úpravy, platné na původním definičním oboru.
\begin{prikladReseniBox}[breakable=false,title={Příklad 1: důkaz roznásobením}]
Ověř pro všechna reálná $x$, že $(x-2)(x^2+2x+4)=x^3-8$.
\tcblower
\[(x-2)(x^2+2x+4)=x^3+2x^2+4x-2x^2-4x-8=x^3-8.\]
Rovnost platí pro každé reálné $x$, protože jsme použili identické úpravy.
\end{prikladReseniBox}
\begin{prikladReseniBox}[breakable=false,title={Příklad 2: pravdivá úprava, ale jiné obory}]
Uprav $\frac{x^2-4}{x-2}$ a posuď, zda lze dosadit $x=2$.
\tcblower
Čitatel $(x-2)(x+2)$ lze krátit \emph{jen při} $x\ne2$. Výsledný zápis $x+2$ je definovaný i pro $x=2$, ale \textbf{původní zlomek ne}. Obecně tedy rovnost platí pouze na $\mathbb R\setminus\{2\}$.
\end{prikladReseniBox}
\begin{prikladReseniBox}[breakable=false,title={Příklad 3: kontrola podílu a zbytku}]
Ověř dělení $x^3+1=(x-1)(x^2+x+1)+2$.
\tcblower
\[(x-1)(x^2+x+1)+2=x^3-1+2=x^3+1.\]
Toto je polynomická identita pro všechna reálná $x$. Podíl číselných hodnot $(x^3+1)/(x-1)$ je však definován jen pro $x\ne1$.
\end{prikladReseniBox}
\blok{Samostatná práce}
\begin{zadaniReseniBox}[breakable=false,title={Procvičení: rozhodni a zdůvodni}]
\begin{enumerate}[label=\alph*)]
\item Ověř $(x+2)(x^2-2x+4)=x^3+8$.
\item Rozlož a zpětně ověř $2x^3-18x$.
\item Uprav $\frac{x^2-9}{x-3}$ včetně podmínek.
\item Najdi protipříklad k $(x-1)^2=x^2-1$.
\item Vyděl $(x^2+1):(x+1)$, určete podíl a zbytek.
\end{enumerate}
\tcblower
\begin{enumerate}[label=\alph*)]
\item Roznásobením $x^3+8$, pro všechna $x\in\mathbb R$.
\item $2x(x-3)(x+3)$, protože $2x[(x-3)(x+3)]=2x(x^2-9)$.
\item $x+3$, původní podmínka $x\ne3$.
\item Pro $x=0$ je vlevo $1$, vpravo $-1$. Správný rozvoj je $x^2-2x+1$.
\item $Q(x)=x-1$, $R(x)=2$, neboť $(x+1)(x-1)+2=x^2+1$.
\end{enumerate}
\end{zadaniReseniBox}
\poznamka{Po testu vybírej příklady podle zjištěných slabin třídy. Nezaměň obecný důkaz s úspěšným dosazením dvou hodnot. Cílem není další objem látky, ale správná argumentace.}
\blok{Volitelná rezerva}
\begin{zadaniReseniBox}[breakable=false,title={Rozšíření: rozpoznání správného mezikroku}]
Žák uvedl $x^3+3x^2-4=x^2(x+3)-4$ a tvrdí, že je hotov rozklad na součin. Je to pravda? Dokonči rozklad.
\tcblower
Rovnost je správná, ale vpravo zůstává \emph{rozdíl}. Platí
\[x^3+3x^2-4=(x-1)(x^2+4x+4)=(x-1)(x+2)^2.\]
Zkouška roznásobením potvrzuje původní mnohočlen.
\end{zadaniReseniBox}
\blok{Domácí procvičení}
\begin{zadaniReseniBox}[breakable=false,title={Domácí úkol: ověřování úprav a podmínek}]
\begin{enumerate}[label=\alph*)]
\item Ověř $(x-1)(x^2+x+1)=x^3-1$.
\item Rozlož $x^2-6x+9$ a ověř roznásobením.
\item Ověř $\frac1{x-1}+\frac1{x+1}=\frac{2x}{x^2-1}$ a uveď podmínky.
\item Ověř $x^3+2x^2-2x-1=(x+1)(x^2+x-3)+2$.
\item Najdi protipříklad k rovnosti $\frac{x^2+1}{x+1}=x$.
\item Uprav $\frac{x^2-1}{x-1}$ a vysvětli význam zakázané hodnoty.
\end{enumerate}
\tcblower
\begin{enumerate}[label=\alph*)]
\item Po roznásobení $x^3-1$; identita pro všechna $x$.
\item $(x-3)^2=x^2-6x+9$.
\item Společný jmenovatel je $(x-1)(x+1)$; v čitateli $(x+1)+(x-1)=2x$. Podmínky $x\ne-1,1$.
\item Roznásobením vznikne $x^3+2x^2-2x-3+2$, tedy levá strana. Identita pro každé $x$.
\item Např. $x=0$: původní zlomek má hodnotu $1$, pravá strana $0$. Původní obor $x\ne-1$.
\item $x+1$, $x\ne1$; upravený výraz existuje také pro $x=1$, původní zlomek nikoli.
\end{enumerate}
\end{zadaniReseniBox}
\end{document}
